\documentclass[../main.tex]{subfiles}

\begin{document}
\section{Method}
\justify
This is the method section. You might want to have tables in the method section, like Table \ref{tab:activityDistribution}.\par\vspace{5mm}

\begin{table}[H]
    \small
    \centering
    \caption{Distribution of the different activities among the different subjects.}
    \label{tab:activityDistribution}
    \begin{tabular}{c c c}
        \hline\hline
        \rowcolor{lightgray}\hspace{6mm}\textbf{Activity}\hspace{6mm} & \hspace{6mm}\textbf{Number Recordings}\hspace{6mm} & \hspace{6mm}\textbf{Number Subjects}\hspace{6mm} \\\hline
        Walking & 6 & 6 \\
        \rowcolor{lightlightgray} Running & 5 & 5\\
        Cycling & 8 & 5\\
       \hline
    \end{tabular}
\end{table}

\noindent
Equations are also frequently used in the method section. You can use equations with numbering as shown in Eq. \ref{eq:normalization}.

\begin{equation}\label{eq:normalization}
    x^{Normalized}_{i} = \frac{x_i-min(x)}{max(x)-min(x)}
\end{equation}

\noindent
Where $x_i$ is each data point in the dataset to be normalized and $x$ is the entire dataset.\par\vspace{5mm}

\noindent
You might want to use equations without numbering, then you do it as shown in Eq. \ref{eq:nonumbering}.

\begin{equation*}\label{eq:nonumbering}
    x_{i} = [PPG\hspace{2mm}Accel_{_{X}}\hspace{2mm}Accel_{_{Y}}\hspace{2mm}Accel_{_{Z}}\hspace{2mm}Gyro_{_{X}}\hspace{2mm}Gyro_{_{Y}}\hspace{2mm}Gyro_{_{Z}}]   
\end{equation*}

\noindent
You can also write math commands within the dollar sign:

\noindent
$x = \frac{2}{4} \cdot \pi$\par


\end{document}
